\section{Restrições e Extensões de Funções}
\label{sec:matematica:funcoes:restricoes-e-extensoes}

\subsection{Restrições}

\begin{defn}[Restrição de uma Função]
	\label{def:matematica:funcoes:restricoes-e-extensoes:restricao}

	Seja

		[
			f:A\to B
		]

	e seja

		[
			X\subseteq A.
		]

	A restrição de (f) a (X) é a função

	[
	f|_X:X\to B
	]

	definida por

		[
			f|_X(x)=f(x)
		]

	para todo (x\in X).

\end{defn}

\begin{obs}

	A restrição preserva a regra de correspondência da função original,
	alterando apenas o seu domínio.

\end{obs}

\begin{prop}
	\label{prop:matematica:funcoes:restricoes-e-extensoes:restricao-funcao}

	A restrição de uma função é uma função.

\end{prop}

\begin{proof}

	Seja (x\in X).

	Como (X\subseteq A) e

		[
			f:A\to B,
		]

	existe um único elemento (f(x)\in B).

	Logo (f|_X) satisfaz as propriedades de existência e unicidade da
	imagem.

\end{proof}

\begin{prop}
	\label{prop:matematica:funcoes:restricoes-e-extensoes:restricao-imagem}

	Seja

		[
			f:A\to B
		]

	e

		[
			X\subseteq A.
		]

	Então

	[
	\operatorname{Im}(f|_X)
		=
		{
			f(x)
			\mid
			x\in X
		}.
	]

\end{prop}

\begin{proof}

	Segue diretamente da definição de imagem.

\end{proof}

\begin{prop}[Restrições Sucessivas]
	\label{prop:matematica:funcoes:restricoes-e-extensoes:restricoes-sucessivas}

	Se

		[
			Y\subseteq X\subseteq A,
		]

	então

	[
	(f|_X)|_Y
	=
	f|_Y.
	]

\end{prop}

\begin{proof}

	Para todo (y\in Y),

	[
			((f|_X)|_Y)(y)
			=
			f|_X(y)
			=
			f(y).
		]

	Pelo princípio da extensionalidade das funções,

	[
			(f|_X)|_Y
			=
			f|_Y.
		]

\end{proof}

\subsection{Extensões}

\begin{defn}[Extensão de uma Função]
	\label{def:matematica:funcoes:restricoes-e-extensoes:extensao}

	Sejam

		[
			f:A\to B
		]

	e

		[
			g:C\to B
		]

	tais que

		[
			A\subseteq C.
		]

	Diz-se que (g) é uma extensão de (f) quando

		[
			g|_A=f.
		]

\end{defn}

\begin{obs}

	Uma extensão preserva completamente o comportamento da função
	original em seu domínio original.

\end{obs}

\begin{ex}

	Considere

	[
	f:\mathbb N\to\mathbb N,
	\qquad
	f(n)=n^2.
	]

	A função

	[
	g:\mathbb Z\to\mathbb Z,
	\qquad
	g(z)=z^2
	]

	é uma extensão de (f).

\end{ex}

\begin{prop}
	\label{prop:matematica:funcoes:restricoes-e-extensoes:extensao-determina}

	Seja

		[
			g:C\to B
		]

	uma extensão de

		[
			f:A\to B.
		]

	Então

	[
	g(a)=f(a)
	]

	para todo (a\in A).

\end{prop}

\begin{proof}

	Como

	[
	g|_A=f,
	]

	segue da definição de restrição que

		[
			g(a)=f(a)
		]

	para todo (a\in A).

\end{proof}

\subsection{Inclusão entre Funções}

\begin{defn}[Inclusão de Funções]
	\label{def:matematica:funcoes:restricoes-e-extensoes:inclusao}

	Sejam

		[
			f:A\to B
		]

	e

		[
			g:C\to D.
		]

	Diz-se que

		[
			f\subseteq g
		]

	quando, considerados como relações,

	[
			f
			\subseteq
			g.
		]

\end{defn}

\begin{prop}
	\label{prop:matematica:funcoes:restricoes-e-extensoes:inclusao-caracterizacao}

	Para funções (f) e (g),

	[
			f\subseteq g
		]

	se, e somente se,

	\begin{enumerate}

		```
		\item
		      \[
			      \operatorname{Dom}(f)
			      \subseteq
			      \operatorname{Dom}(g);
		      \]

		\item para todo
		      \(x\in\operatorname{Dom}(f)\),

		      \[
			      f(x)=g(x).
		      \]
		      ```

	\end{enumerate}

\end{prop}

\begin{proof}

	Segue imediatamente da definição de inclusão entre relações e da
	definição de função.

\end{proof}

\begin{prop}
	\label{prop:matematica:funcoes:restricoes-e-extensoes:ordem-parcial}

	A relação de inclusão entre funções é uma ordem parcial.

\end{prop}

\begin{proof}

	A inclusão entre relações é reflexiva, antissimétrica e transitiva.

	Como funções são relações, as mesmas propriedades permanecem válidas.

\end{proof}

\begin{prop}
	\label{prop:matematica:funcoes:restricoes-e-extensoes:restricao-inclusao}

	Seja

		[
			f:A\to B
		]

	e

		[
			X\subseteq A.
		]

	Então

	[
	f|_X
	\subseteq
	f.
	]

\end{prop}

\begin{proof}

	Para todo (x\in X),

	[
			f|_X(x)
			=
			f(x).
		]

	Logo o gráfico de (f|_X) é subconjunto do gráfico de (f).

\end{proof}

\begin{prop}
	\label{prop:matematica:funcoes:restricoes-e-extensoes:extensao-inclusao}

	Seja

		[
			g
		]

	uma extensão de

		[
			f.
		]

	Então

	[
	f
	\subseteq
	g.
	]

\end{prop}

\begin{proof}

	Como

	[
	g|_{\operatorname{Dom}(f)}
	=
	f,
	]

	todo par ordenado pertencente a (f) também pertence a (g).

	Portanto

		[
			f\subseteq g.
		]

\end{proof}

\end{document}
